% Created 2025-11-04 Tue 15:39
% Intended LaTeX compiler: pdflatex
\documentclass[11pt]{article}
\usepackage[utf8]{inputenc}
\usepackage[T1]{fontenc}
\usepackage{graphicx}
\usepackage{longtable}
\usepackage{wrapfig}
\usepackage{rotating}
\usepackage[normalem]{ulem}
\usepackage{amsmath}
\usepackage{amssymb}
\usepackage{capt-of}
\usepackage{hyperref}
\usepackage{amsmath} \usepackage{unicode-math} \usepackage{cancel} \usepackage{graphicx}
\renewcommand{\theta}{\vartheta} \renewcommand{\phi}{\varphi} \renewcommand{\epsilon}{\varepsilon}
\newcommand{\odv}[1]{\dot{\vec{#1}}} \newcommand{\oddv}[1]{\ddot{\vec{#1}}}
\newcommand{\rot}{\mathrm{rot}}\newcommand{\dive}{\mathrm{div}} \newcommand{\iu}{\mathrm{i}}
\newcommand\mapsfrom{\mathrel{\reflectbox{\ensuremath{\mapsto}}}}
\date{\today}
\title{Vprasanja}
\hypersetup{
 pdfauthor={},
 pdftitle={Vprasanja},
 pdfkeywords={},
 pdfsubject={},
 pdfcreator={Emacs 30.2 (Org mode 9.7.31)}, 
 pdflang={English}}
\begin{document}

\maketitle
\tableofcontents

\section{Kolokvij 1, 1999/2000}
\label{sec:org072f85e}

V rešitvah po uporabi nastavka \(u(x, t) = z(x) e^{- \mathrm{i} \omega t}\) dobi enačbo

\[ \frac{\mathrm{d} ^2 z }{\mathrm{d} x ^2} + k ^2 z = 0.
\]

Nastavek, ki naj bi rešil to enačbo, je potem zapisan kot

\[ z (x) = \sin (kx) + \alpha \cos (kx).
\]

Moje vprašanje je, zakaj ima sinus koeficient 1 in kosinus poljuben koeficient? Nisem našel načina, kako bi iz robnih pogojev prišel do takega spoznanja. So ti koeficienti posledica fizikalne interpretacije?
\section{Kolokvij 1, 2000/2001}
\label{sec:org3be468d}
\subsection{1. vprašanje}
\label{sec:org2b0c815}
Valovni vektor \(k_2\) ima v rešitvah vrednosti \(k_2 = \sqrt{2} k_1\), ko se valovanje loči na razpleteni struni. Zakaj?

Razumem, da \(\sqrt{2}\) pride iz \(\sin / \cos \frac{\pi}{4}\), vendar kam izgine \(\frac{1}{2}\)?

Je mogoče to posledica tega, da je celoten valovni vektor v \(x\)-smeri, ki je usmerjena vzporedno z valovanjem, vsota posameznih dveh? Torej

\[ k_{2x} = 2 \cdot \frac{\sqrt{2}}{2} k_1?
\]
\subsection{2. vprašanje}
\label{sec:orgaca966d}

Ko iščemo oddaljenost \(a\) in maso \(m\), ki se nahaja pri \(x = 0\), da se val ne odbije, je drugi Newtonov zakon za utež

\begin{equation}
\label{eq:1}
 \left. m \ddot{u}_{pred} \right|_{x = 0} = - F_0 \left. \frac{\partial u_{pred}}{\partial x} \right|_{x = 0} + F_0 \left. \frac{\partial u_{po}}{\partial x} \right|_{x = 0},
\end{equation}

kjer je

\[ u(x, t) = \begin{cases}
              u_{pred} (x, t) = e^{\mathrm{i} (k_1 x - \omega t)} + R e^{\mathrm{i} (- k_1 x - \omega t)} &; x < 0 \\
              u_{po} (x, t) = T_1e^{\mathrm{i} (k_1 x - \omega t)} + R_1 e^{\mathrm{i} (-k_1 x - \omega t)} & ; a < x < 0 \\
              u_{Y} (x, t) = Te^{\mathrm{i} (k_2 x - \omega t)} &; a < x
             \end{cases}
\]

Torej \(u_{pred}\) je pred maso, \(u_{po}\) je po masi in \(u_Y\) je valovanje na razpletenih strunah.

Newtonov zakon potem po upoštevanju definicij, odvajanju in \(x = 0\) postane

\[ - \omega ^2 m (1 + R) = \mathrm{i} k_1 F_0 \left( -1 + R + T_1 - R_1 \right).
\]

V rešitvah prvega faktorja \((1 + R)\) ni. So rešitve napačne ali jaz kje počnem matematično katastrofo?

Ali je pri Newtonovem zakonu \ref{eq:1} važno, kateri nastavek za rešitev uporabimo pri masi? Bosta zapis \ref{eq:1} ter

\[  \left. m \ddot{u}_{po} \right|_{x = 0} = - F_0 \left. \frac{\partial u_{pred}}{\partial x} \right|_{x = 0} + F_0 \left. \frac{\partial u_{po}}{\partial x} \right|_{x = 0},
\]

pripeljala do istega rezultata?
\end{document}
